\subsection{Galois coverings}

DEFINITION 3. --- Let B be a nonempty topological space. A covering E of B is said to be Galois if it is connected and if, for every point b of B, the action of the group \(\mathrm{Aut}_{B}(E)\) of B-automorphisms of E on the fiber \(E_{b}\) is transitive.

Let E be a Galois covering of a topological space B and let p be its projection. The map p is surjective (A, I, p. 56, def. 6), so the space B is connected. Consequently, the covering (E, p) has a degree, which is nonzero.

PROPOSITION 5. --- Let B be a nonempty topological space and let E be a Galois covering of B. The map \((h,x)\mapsto h(x)\) from \(\mathrm{Aut}_{\mathrm{B}}(\mathrm{E})\times\mathrm{E}\) into E is a right action of the group \(\mathrm{Aut}_{\mathrm{B}}(\mathrm{E})^{\circ}\) on E which makes E a principal covering with group \(\mathrm{Aut}_{\mathrm{B}}(\mathrm{E})^{\circ}\) .

Endow the group \(\mathrm{Aut}_{\mathrm{B}}(\mathrm{E})^{\circ}\) with the discrete topology; the operation law \((h,x)\mapsto h(x)\) is then continuous. This action is free (I, p. 34, corollary 2 of proposition 11). Since E is a Galois covering of B, its fibers are the orbits of this action. By proposition 4 of I, p. 97, E is a principal covering with group \(\mathrm{Aut}_{\mathrm{B}}(\mathrm{E})^{\circ}\) .

THEOREM 2. --- Let B be a connected topological space, and let E be a connected, nonempty covering of B. The following properties are equivalent:

\begin{enumerate}
\def\labelenumi{(\roman{enumi})}
\item
  The covering E is Galois.
\item
  There exists a discrete topological group G and a continuous right action of G on E that makes E a principal covering with group G.
\item
  The covering \(E \times_{B} E\) of E defined by the projection \(\mathrm{pr}_{1}\) is trivializable.
\end{enumerate}

When these conditions are satisfied, the canonical map \(G \to \mathrm{Aut}_{B}(E)^{\circ}\) defined by the action of G is an isomorphism of groups.

If, moreover, the space B is locally connected, the preceding properties are equivalent to the following property:

(i') There exists a point b of B such that the action of the group \(\mathrm{Aut}_{\mathrm{B}}(\mathrm{E})\) on the fiber \(E_{b}\) is transitive.

The implication (i)⇒(ii) follows from Proposition 5, and the implication (ii)⇒(iii) from Remark 1 of I, p. 95. Let us prove (iii)⇒(i). Denote by \(p\colon E \to B\) the projection of the covering E. Since B is connected, this covering has a degree, and this degree is nonzero because E is not empty. Let \(b\) be a point of B. Let us prove that the action of \(\mathrm{Aut}_{B}(E)\) on the fiber \(E_{b}\) is transitive. By the preceding, this fiber is nonempty. Let \(x\) and \(x'\) be points of \(E_{b}\). The B-morphisms \(h\colon E \to E\) such that \(h(x) = x'\) correspond bijectively to continuous sections \(s\) of the map \(\mathrm{pr}_{1}\colon E \times_{B} E \to E\) such that \(s(x) = (x, x')\) (I, p. 9, prop. 3). Under hypothesis (iii), such a section exists (I, p. 70, cor. 2) and is unique (I, p. 34, cor. 1 of prop. 11) because the space E is connected. Thus, for every pair \((x, x') \in E \times_{B} E\) there exists a unique B-morphism \(h\colon E \to E\) such that \(h(x) = x'\) . If \(h'\colon E \to E\) is the unique B-morphism such that \(h'(x') = x\) , we have \(h'(h(x)) = x\) and \(h(h'(x')) = x'\) , whence \(h' \circ h = \text{Id}_E\) and \(h \circ h' = \text{Id}_E\) . This proves that \(h\) is a B-automorphism of E. The covering E is therefore Galois.

We have shown that conditions (i), (ii), (iii) are equivalent. Suppose they are satisfied. Let \(\delta\colon\mathbf{G}\to\mathrm{Aut}_{\mathrm{B}}(\mathrm{E})\) be the map which associates with \(g\in\mathrm{G}\) the B-automorphism \(x\mapsto x\cdot g\) of E. The map \(\delta\) is a group homomorphism from G into \(\mathrm{Aut}_{\mathrm{B}}(\mathrm{E})^{\circ}\). Since G acts freely on E, the map \(\delta\) is injective. Let \(h\in\mathrm{Aut}_{\mathrm{B}}(\mathrm{E})\), let \(x\in\mathrm{E}\), and let \(g\) be the unique element of G such that \(x\cdot g=h(x)\).

B-morphisms \(h\) and \(\delta(g)\) coincide at \(x\) and therefore everywhere, since the space E is connected (I, p. 34, cor. 1 of prop. 11), which proves that \(\delta\) is an isomorphism.

Suppose now that the space B is locally connected, and let us prove, under hypothesis (i'), that the covering E is Galois. Let \(E'\) be the quotient space of E by the right action of \(\mathrm{Aut}_{B}(E)^{\circ}\) and denote by \(q\colon E \to E'\) the canonical map. It makes E a surjective covering of E' (I, p. 98, example 2); the map \(p\colon E \to B\) passes to the quotient to define a continuous map \(p'\colon E' \to B\) such that \(p'\circ q = p\) . Since the space B is locally connected, the B-space \((E', p')\) is a covering (I, p. 81, prop. 7). Under hypothesis (i'), there exists a point \(b\) of B such that the fiber \(E'_{b}\) has exactly one element; since the space B is connected, the same holds for every point \(b\) of B (I, p. 74, prop. 4), which proves that E is a Galois covering of B.

PROPOSITION 6. --- Let B be a topological space and G a group. Let (E, p) be a principal covering of B with group G. Suppose the space B is connected and locally connected. Let E₀ be a connected component of E and let G₀ be the subgroup of G stabilizing E₀ (A, I, p. 51). The B-space (E₀, p \textbar{} E₀) is a principal covering for the right action of G₀ on E₀. This covering is Galois.

Since the space B is locally connected, so is E, so that \(E_{0}\) is an open subspace of E (TG, I, p. 85, prop. 11). Since it is also closed (TG, I, p. 83), the space \(E_{0}\) is therefore a covering of B (I, p. 80, cor. 1). Since B is connected, this covering has a degree; since \(E_{0}\) is nonempty, all fibers of \(p|E_{0}\) are therefore nonempty. Let x be a point of \(E_{0}\), and let \(x' \in E_{0}\) such that \(p(x') = p(x)\); there exists an element \(g \in G\) such that \(x \cdot g = x'\). Since the connected components \(E_{0}\) and \(E_{0} \cdot g\) have a common point, they are equal, which implies \(g \in G_{0}\). Thus, the group \(G_{0}\) acts transitively on the fiber of the B-space \(E_{0}\) at \(p(x)\). Proposition 6 follows.

Note. --- If, in Proposition 6, one does not assume the space B to be locally connected, it may happen that the space \(E_{0}\) is not a covering of B (I, p. 147, exerc. 1).

